\scanpage{454}

formulas which have a meaning only if we suppose given embeddings (or at least morphisms)

\diagram{
\[
(155)\quad
\left\{
\begin{tikzcd}[ampersand replacement=\&,column sep=large]
T_\rho \arrow[r,hook] \& G \\
T_\sigma \arrow[r,hook] \& G
\end{tikzcd}
\right.
\qquad
\begin{gathered}
\text{and a subgroup}\\
\begin{tikzcd}[ampersand replacement=\&,column sep=large]
T_0 \arrow[r,hook] \& G
\end{tikzcd}
\end{gathered}.
\]
}

When we have such embeddings, the formulas (154) have a meaning, and define \deleted{\unsure{all the}} (with the quantities (152) known) the quantities (153) in terms of one \deleted{\unsure{of them}}\inserted{of them}, say $\underline u$, or \deleted{\illegible{1}} $f$, or $\alpha$, or $\beta$ (but not in terms of $\alpha_0$ or of $\beta_0$).

In fact, we have a subgroup

\diagram{
\[
\begin{array}{r@{\;}c@{\;}l}
\begin{gathered}(156)\\[2em]\deleted{\text{\unsure{such}}}\\[3em]{}\end{gathered}
&\begin{tikzcd}[ampersand replacement=\&,row sep=3em]
{}\arrow[dd]\\{}\\{}
\end{tikzcd}
&\begin{gathered}
T_0\subset G\\
\text{such that}\quad
\left\{\begin{array}{l}
\begin{tikzcd}[ampersand replacement=\&,column sep=large]
T_0 \arrow[r,"{\substack{\delta_G|T_0\\\sim}}"] \& G_m
\end{tikzcd}\\[5pt]
\text{i.e.}\quad G\simeq T_0.SG\\
\hspace{6em}\text{product }1/2\text{ direct},
\end{array}\right.
\end{gathered}
\end{array}
\]
}

so that \deleted{the quantity}\inserted{for $U\in\Gamma G$,} there exists a \uline{unique} \deleted{$f\in\Gamma SG$ such that}\inserted{pair $(f,u)\in\Gamma(SG\times T_0)$ with}
\[
U=f^{-1}\underline u.
\]
hence unique $\alpha,\beta,\alpha_0,\beta_0\in\Gamma G$ satisfying (with the preceding $f,u$) (154) :
\[
(157)\quad
\alpha=\underline u\,\underline a^{-1},\quad
\beta=\underline u\,\underline b^{-1},\quad
\alpha_0=f^{-1}\alpha=f^{-1}\underline u\,\underline a^{-1},\quad
\beta_0=f^{-1}\deleted{\unsure{\alpha}}\inserted{\beta}=f^{-1}\underline u\,\underline b^{-1}.
\]
\mspage{453}

\scanpage{455}
\mspage{530}

We will moreover have
\[
\deleted{\unsure{\det\underline U=}}\det\underline u
=\det\underline U=\mu
\]
\[
\begin{aligned}
\det\alpha&=\underline{\det}(\underline u)\det(\underline a)^{-1}=\varepsilon'
&&\text{by (152')}\\
\det\beta&=\det(\underline u)\det(\underline b)^{-1}=1\\[3pt]
\text{and}\quad\det\alpha_0&=\det\alpha=\varepsilon'\\
\det\beta_0&=\det\beta=1
\end{aligned}
\]
we will thus indeed have
\[
\begin{aligned}
\alpha,\alpha_0&\in\Gamma G^{\pm1}\overset{\text{definition}}=
\{g\mid(\det g)^2=1\}\\
\beta,\beta_0&\in\Gamma SG
\end{aligned}
\]
--- but we have assumed, of course, that the \unsure{homomorphisms} $\delta_\rho,\delta_\sigma$ of (135) are \unsure{induced}, via the embeddings (155), by the determinant $\delta_G$.

To arrive at the relation \unsure{or} \unsure{of equation}
\[
(157)\qquad
\underbrace{[\alpha,\rho]}_{\xi}
=\varepsilon\underbrace{[\beta,\sigma]}_{\eta}\varepsilon_1(\mu-1),
\]
we must give ourselves sections $\rho,\sigma$ of $G$ (\deleted{\unsure{or}} more precisely probably, \inserted{resp.\textsuperscript{t}} of maximal $T_\rho$ and $T_\sigma$), and a strict unipotent $\varepsilon_1\in\Gamma SG$, for the relation (157) to make sense.
Let us replace $\alpha,\beta$ by \unsure{some|the} values $\underline u\,\underline a^{-1},\underline u\,\underline b^{-1}$, the relation becomes (using the commutativity of $T_\rho$, \deleted{\unsure{and}} which we had not really brought \unsure{into play} until now)
\[
(*)\qquad
[u,\rho]
=\varepsilon\cdot\underline u
\bigl(\underline b^{-1}\sigma(\underline b)^{\deleted{\unsure{-1}}}\bigr)^{\deleted{\unsure{-1}}}
\sigma(\underline u)^{-1}\varepsilon_1(\mu-1).
\]
\mspage{454}

\scanpage{456}

and will follow from the two conditions together
\[
(158)\quad
\left\{
\begin{aligned}
\text{a)}\quad&b^{-1}\sigma(b)=\varepsilon_\sigma(b)
&&\text{\uline{identically} if }b\in\Gamma N'_\sigma\\
&\text{i.e.}\quad b(\sigma)=\varepsilon_\sigma(b).\sigma
&&\text{identically}\\[4pt]
\text{b)}\quad&[u,\rho]=[u,\sigma]\varepsilon_1(\mu-1)
&&\text{identically in }u\in\Gamma T_0\\
&&&\text{for }\mu=\delta_{T_0}(u)\overset{\mathrm{dfn}}=\det(u)\\
&\text{i.e.}\quad\rho(u)=\varepsilon_1(1-\mu)\sigma(u)
&&\text{identically in }u
\end{aligned}
\right.
\]
which can also be written (applying $\sigma^{-1}$ to the two \unsure{sides|sides})
\[
\sigma^{-1}\rho(u)=\varepsilon_0(1-\mu).u
\qquad\text{where}\quad\varepsilon_0=\sigma^{-1}(\varepsilon_1)
\]
and is hardly conceivable unless \uline{$T_0$ normalizes $L_0$}, and unless $\sigma^{-1}\rho\in\Gamma T_0.L_0$ \deleted{and} is of the form
\[
(159)\quad
\left\{
\begin{aligned}
\text{a)}\quad&\widetilde{\sigma^{-1}\rho}=\widetilde{\varepsilon_0u_0}
\qquad u_0\in\Gamma T_0\\
\text{b)}\quad& T_0\text{ operates on }L_0\text{ by the \unsure{coaction|action} \unsure{induced}}\\
&\hspace{8em}\text{by the determinant}
\end{aligned}
\right.
\]
(There is all the \unsure{same} still some leeway, \unsure{in relation|by} to the draconian conditions envisaged in \uline{$\overline{\mathrm{VII}}$} !)

But here equation (157) came a bit like hair in soup --- there must be more convincing and demanding ways of bringing it in namely \deleted{by} in connection with formulas
\[
(160)\quad
\left\{
\begin{aligned}
\underline U(\rho)&=\operatorname{int}(\alpha_0).\rho
&&\bigl(=\alpha_0\rho(\alpha_0)^{-1}\rho\bigr)\\
\underline U(\deleted{\unsure{\rho}}\inserted{\sigma})&=
\varepsilon\operatorname{int}(\beta_0)\sigma
&&\bigl(=\beta_0\sigma(\beta_0)^{-1}\bigr)\\
\underline U(\varepsilon_0)&=\operatorname{int}(f^{-1})
\bigl(\varepsilon_0^{\deleted{\unsure{\mu}}}\unsure{(\mu)}\bigr)
\end{aligned}
\right.
\qquad
\begin{aligned}
\text{i.e.}\quad[\underline U,\rho]&=\alpha\rho(\alpha)^{-1}\\
[\underline U,\sigma]&=\varepsilon\beta\sigma(\beta)^{-1}
\end{aligned}
\]
\margin{watch the sign $\varepsilon$}

The third relation follows at once from $\underline U=f^{-1}\underline u$, and from the fact that by hypothesis \inserted{(159 b)} $\underline u\in\Gamma T_0$ operates on $L_0$ (\unsure{identically}) with the multiplier $\mu=\det(\underline u)$.
\mspage{455}

\scanpage{457}
\mspage{531}

The other two formulas are \unsure{made explicit} successively, by writing $\alpha=\underline u\,\underline a^{-1}$, $\beta=\underline u\,\underline b^{-1}$, $\underline U=f^{-1}\underline u$
\[
\alpha_0=f^{-1}\alpha,\quad\beta_0=f^{-1}\beta
\]
\[
(161)\qquad
\left\{
\begin{aligned}
\underline u(\rho)&=\operatorname{int}(\alpha)(\rho)\\
\underline u(\sigma)&=\varepsilon\operatorname{int}(\beta)(\sigma)
\end{aligned}
\right.
\qquad\text{where }f\text{ no longer appears}
\]
\[
\begin{aligned}
\text{i.e.}\qquad\underline u(\rho)&=\operatorname{int}(\underline u)(\rho)
&&\text{since }\operatorname{int}(\underline a^{-1})\rho=\rho\\
&&&\text{(as }T_\rho\text{ commutative)}\\
&\text{OK !}\\
\underline u(\sigma)&=\deleted{\illegible{1}}\inserted{\varepsilon}\operatorname{int}(u)
\underbrace{\bigl(\underline b^{-1}(\sigma)\bigr)}_{\substack{=\varepsilon.\sigma\\\text{cf. (158)(a)}}}
=\operatorname{int}(u)\sigma\quad\text{OK !}
\end{aligned}
\]

To prove (160) we have used only (158 a) and (159) (b). Let us now try to deduce from them the ``loop relation'' (157) --- it must be a formal calculation, which I must already have done a dozen times, in contexts which become \unsure{differentiated} as we go along ! We eliminate $f$ by writing the relations (160) in the equivalent form (161) and
\[
(161\ \text{bis})\quad
\underline u(\varepsilon_0)=\varepsilon_0(\mu)
\qquad\text{i.e.}\quad
\underbrace{[\underline u,\varepsilon_0]}_{u\varepsilon_0(u)^{-1}}
=\varepsilon_0(\mu-1)\qquad\deleted{\unsure{\operatorname{int}}}
\]
\[
\deleted{\text{i.e.}\quad u\varepsilon_0(u)^{-1}}
\]
\[
\begin{array}{l|ll}
\begin{aligned}[t]
\text{now}\qquad&\\
[\underline u,\varepsilon_0]&=
u\underbrace{\varepsilon_0(u)^{-1}}_{\substack{\unsure{\text{i.e.}}\ \sigma^{-1}\rho(u)^{\deleted{\unsure{-1}}}}}\\
\text{i.e.}\quad\sigma([u,\varepsilon_0])&=\sigma(u)\rho(u)^{-1}
\end{aligned}
&
\begin{aligned}[t]
\text{now}\quad&\deleted{\rho(u)=\rho u\rho^{-1}=}\\
&u\rho u^{-1}\rho^{-1}=\xi\\
&\rho(u^{-1})=u^{-1}\xi\\
&\rho(u)=\xi^{-1}u\\
\unsure{\text{of}=|\text{likewise}}\quad&\sigma(u)=\deleted{\unsure{\varepsilon}}\eta^{-1}u
\end{aligned}
&
\begin{aligned}[t]
\xi&\overset{\mathrm{dfn}}{=}\alpha\rho(\alpha)^{-1}\\[4em]
\eta&\overset{\mathrm{dfn}}{=}\varepsilon\beta\sigma(\beta)^{-1}
\end{aligned}
\end{array}
\]
\[
\text{hence}\qquad
\sigma([u,\varepsilon_0])=
\deleted{\unsure{\xi^{-1}u\eta^{-1}u}}\eta^{-1}\xi
\]
and \unsure{(*)|($\alpha$)} becomes, on applying $\unsure{\sigma}$
\[
\eta^{-1}\xi=\varepsilon_1(\mu-1)
\qquad\text{i.e.}\quad\xi=\eta\varepsilon_1(\mu-1)
\qquad\text{i.e. (157)}
\]
\hfill OK !
\mspage{456}

\scanpage{458}

In summary :

\uline{Proposition} Let us suppose given, in addition to the data of the preceding section $7^\circ$ (page 526' and following \unsure{\}|)}, in particular (135)(136)), \deleted{sections $\rho$ of $\unsure{T_{\rho,\sigma}|T_\rho,\sigma}$} the following

\diagram{
\[
(162)\quad\left\{
\begin{array}{ll}
\text{a)}&\text{section }\rho\in\Gamma T_\rho,\quad\sigma\in\Gamma T_\sigma\quad\text{\uline{central} sections}\\[4pt]
\text{b)}&\text{Embeddings}\quad
\begin{tikzcd}[ampersand replacement=\&,column sep=large]
T_\rho \arrow[r,hook] \& G
\end{tikzcd},\quad
\begin{tikzcd}[ampersand replacement=\&,column sep=large]
N'_\sigma \arrow[r] \& G
\end{tikzcd}\\[4pt]
\text{c)}&\text{torus}\quad
\begin{tikzcd}[ampersand replacement=\&,column sep=large]
T_0 \arrow[r,hook] \& G
\end{tikzcd}\quad\text{such that }\delta_G:\quad
\begin{tikzcd}[ampersand replacement=\&,column sep=large]
\unsure{T_0} \arrow[r,"\sim"] \& G_m
\end{tikzcd}\\
&\text{i.e.}\quad G\simeq T_0.SG\quad\text{(semidirect)}\\[4pt]
\text{d)}&\varepsilon_0\in\Gamma G\quad\text{\deleted{strict} regular unipotent, defining}\\
&\text{thus}\quad
\begin{tikzcd}[ampersand replacement=\&,column sep=large]
\unsure{\lambda} \arrow[r,mapsto] \& \varepsilon_0(\lambda)
\end{tikzcd}:\quad
\begin{tikzcd}[ampersand replacement=\&,column sep=large]
G_a \arrow[r,hook] \& SG
\end{tikzcd}
\end{array}\right.
\]
}

and let us suppose \deleted{(in addition to the \uline{commutativity} of $T_\rho,T_\sigma$}
\[
\begin{aligned}
\inserted{\text{a)}\quad\delta_G|T_\rho=\delta_\rho,\quad\delta_G|N'_\sigma=\delta_\sigma}\\
\deleted{\text{a)}\quad T_0\text{ \unsure{normalizes} }\unsure{L_0}}
\end{aligned}
\]
\[
(163)\quad
\begin{aligned}
\text{b)}\quad&\operatorname{int}(\underline b)(\sigma)=\varepsilon_\sigma(\underline b)\sigma
\quad\text{in }G\text{ identically}\\
&\hspace{12em}\text{for }\underline b\in\Gamma\unsure{N'_\sigma|N^2_\sigma}\\[4pt]
\text{c)}\quad&T_0\text{ normalizes }L_0,\text{ and we have}\\
&u(\varepsilon_0(\lambda))=
\varepsilon_0\bigl(\deleted{\unsure{\delta_G(u)}}\mu\lambda\bigr)
\quad\text{where }\mu=\delta_G(u)\\[4pt]
\text{d)}\quad&\widetilde{\unsure{\sigma|\rho}}^{-1}\widetilde\rho
\in\deleted{\unsure{\widetilde T_0}}\Gamma\widetilde\varepsilon_0\widetilde L_0
\quad\text{i.e. }\exists\text{locally}^{\text{t}}\text{ section }\lambda_0\in\Gamma G_m\\
&\text{with}\quad\lambda_0\widetilde\sigma^{-1}\widetilde\rho=\varepsilon_0u_0,
\quad\text{with }u_0\in\Gamma T_0
\end{aligned}
\]
\margin{\unsure{links} $L_0,T_0$}
\margin{$\begin{array}{l|}\text{\unsure{links}}\quad L_0,T_0\\\rho,\sigma\end{array}$}

Under these conditions, we have for every system (152) satisfying the relations (152') (i.e. in fact for every \deleted{\unsure{syst}} section of $\mathcal M_{\rho,\sigma}$ !) a unique \unsure{system} (153), satisfying the relations (154) --- namely $\underline U$ can be written uniquely in the form
\[
(164)\quad\underline U=f^{-1}\underline u,\qquad f\in\Gamma SG,\quad\underline u\in\Gamma T_0
\]
and we put
\[
\begin{aligned}
(165)\quad\alpha&=\underline u\,\underline a^{-1},\quad
\beta=\underline u\,\underline b^{-1},\quad
\alpha_0=f^{-1}\unsure{\alpha|\alpha_\rho}=
\deleted{\unsure{\beta}}\inserted{f}^{-1}\underline u\,\underline a^{-1},\\
\beta_0&=f^{-1}\beta=f^{-1}\underline u\,\underline b^{-1}.
\end{aligned}
\]
\mspage{457}

\scanpage{459}
\mspage{532}

All this being laid down, we then have the formulas
(160), (161), (161 bis), and the ``loop'' formula (157).

We now ask ourselves the question
under what conditions knowledge of
$(\alpha,\beta,f)$ (for example) among the quantities (153),
allows us to recover the quantities (152) ---
i.e. a section of $\mathcal M_{\rho,\sigma}$ --- and what conditions
must be imposed on $(\alpha,\beta,f)$. It
would suffice \inserted{(for uniqueness)} to be able to construct $\underline u$, whence
$\underline U,\underline a,\underline b$ by the first three formulas
(154). \deleted{\unsure{If, on the other hand we have}} We must
therefore find $\underline a\in\Gamma T_\rho$ such that $\alpha\underline a$ is a section
of $T_0$, --- indeterminacy in $\underline a$ is by
$\Gamma T_0\cap T_\rho$, and $\underline b\in\Gamma N'_\sigma$ such that $\beta\underline b$ is a section
of $T_0$ --- indeterminacy in \deleted{$\unsure{\Gamma N'_\sigma\cap}$} $\Gamma T_0\cap N'_\sigma$.
The indeterminacy in $\underline u$ is right
multiplication by a section of $T_0\cap T_\rho\cap N'_\sigma$, so it
will suffice that we have
\[
\text{(164)}\qquad T_0\cap T_\rho\cap N'_\sigma=\{1\}.
\]
\margin{in fact, we will \unsure{even} have $T_0\cap T_\rho=\{1\}$ in the cases \unsure{considered} in \uline{II}, \uline{III} and \unsure{following|sequel}\ldots}
Assuming this relation satisfied, what
conditions to put on $(\alpha,\beta,f)$ to come
from a \unsure{system} (152) --- i.e. from a section of $\mathcal M_{\rho,\sigma}$ ?

\mspage{458}

\scanpage{460}

We see that \unsure{if} a section \deleted{$\unsure{\varphi}$} $\varphi_0$ of $\mathcal M_{\rho,\sigma}$, \deleted{is}
\deleted{\unsure{replaced}} corresponding to $\alpha_0,\beta_0,f_0$, is replaced
by $\varphi_1\overset{\unsure{\text{definition}}}{=}f^{-1}.\varphi_0$, corresponding to $\alpha_1,\beta_1,f_1$, then
\[
\text{(165)}\qquad \alpha_1=\alpha_0,\quad\beta_1=\beta_0,\quad f_1=f_0f,
\]
for we see that if $\varphi_0$ corresponds to $(\underline U_0,\underline a_0,\underline b_0)$,
then $\varphi_1$ to $(U_1,\underline a_1,\underline b_1)$, then we have
\[
\text{(166)}\qquad U_1=f^{-1}U_0,\quad\underline a_1=\underline a_0,\quad\underline b_1=\underline b_0
\]
(for $\Theta_G\mid SG=$ the canonical inclusion $SG\hookrightarrow G$,
and $\Theta_\rho$\inserted{$,\Theta_\sigma$} \deleted{is} trivial on $SG$). \deleted{Thus This
shows that} \inserted{Thus \deleted{\unsure{it}}} we will have
\[
u_1=u_0\qquad (\text{for }\det U_1=\det U_0)
\]
whence (165). This shows that the conditions
sought on $(\alpha,\beta,f)$ depend only on
$\alpha,\beta$, not on $f$. We would like them
to be expressed by the validity of the loop
formula\inserted{(157)}, for a parameter \inserted{$\varepsilon,$} $\mu$ not specified $(\varepsilon^2=1)$
(which \unsure{will be} then uniquely determined --- and \unsure{will allow}
the \unsure{secondary definition} of \unsure{local notion|local action}\ldots).

\margin{to tell the truth, this condition is not \uline{sufficient} even in the cases considered --- we will find \unsure{therefore|in} a $\underline u\in B'_0$, and we must require that $\underline u\in T_0$ (already for $(\alpha,\beta)$ to be ``normalized'') i.e. $B=0$, where $B=ac+bd+bc$, if $\underline u=\left(\begin{smallmatrix}a&b\\c&d\end{smallmatrix}\right)$ (modular case) cf. proposition page~412'\ldots\ \unsure{To have it} $\unsure{R}$ ; \unsure{itself being} \unsure{additional condition}, \unsure{it|which} \unsure{is necessary to replace|is necessary to consider} $\mathcal M_{\rho,\sigma}$ \unsure{by the product} $\tfrac12$ \unsure{direct by} $L_0$, \unsure{on which we set} $\unsure{f_\Delta}=\unsure{0}$, \unsure{i.e.} $S\mathcal M_{\rho,\sigma}$.}

The sufficiency of this condition \unsure{will be}
the subject of the reflections in (II) (III),
\unsure{under} conditions (more draconian than those
\unsure{of the} preceding proposition (page~531') which we
have imposed on the triple $(\varepsilon_0,\rho,\sigma)$. It

\mspage{459}

\scanpage{461}
\mspage{533}

is possible that the conditions at which we
\inserted{\unsure{there}} arrived were more restrictive
than necessary, and that I shall be obliged
to come back to this later, in a context
more general than that of forms of
$\mathrm{Gl}(2)$, $\mathrm{GP}(1)$, $\mathrm{Sl}(2)$, but for the moment
the detailed study made by taking our inspiration
from the hypotheses brought out in \deleted{\unsure{the}} II, III will
have served to clear the ground for us, to clarify
our ideas on an important particular
case.

\uline{Remark.} Let us recall that we have seen that in
fact $\mu$ is already defined in terms of $\alpha$ alone,
and even in terms of $\alpha\rho(\alpha)^{-1}=\xi$ \unsure{;} \deleted{\unsure{indeed|therefore}}
But it is determined in terms of $\beta$
--- in fact, of $\beta\sigma(\beta)^{-1}$ --- only modulo
knowledge of $\varepsilon$ (in practice, modulo
sign\ldots). But if we impose (for example) $\varepsilon=1$,
the datum of $\alpha$ determines $\beta$ only
modulo \deleted{$\unsure{\beta}$} $\beta\longrightarrow\beta T_\sigma$\ldots

\mspage{460}

\scanpage{462}

The structure of $\mathcal M_{\rho,\sigma}$ being well understood,
let us return to the group $S\mathcal M_{\rho,\sigma}$\inserted{(cf. page~495')}, considered as
the semidirect product of $\mathcal M_{\rho,\sigma}$ by the subgroup
$L_0$ of $SG$, on which $\mathcal M_{\rho,\sigma}$ in fact operates
via $\chi$. Denoting by

\noindent (165)\quad
\diagram{\begin{tikzcd}[ampersand replacement=\&,column sep=large]
\deleted{\unsure{\varepsilon_0}:\unsure{U}}\;\inserted{\unsure{\mathcal E|\varepsilon}_{L_0}:G_a} \arrow[r,"\sim"] \&
\operatorname{Ker}\bigl(S\mathcal M_{\rho,\sigma}\longrightarrow\mathcal M_{\rho,\sigma}\bigr)
\end{tikzcd}}

\inserted{corresponding to $\varepsilon_0$,}
we will have by \deleted{\unsure{definition}} construction
\[
\text{(166)}\qquad
U\bigl(\unsure{\mathcal E|\varepsilon}_{L_0}(\lambda)\bigr)
=\unsure{\mathcal E|\varepsilon}_{L_0}(\mu\lambda)
\qquad\text{where }\mu=\chi(U).
\]

\deleted{The group $S$}

\mspage{461}

\scanpage{463}
\mspage{534}

\deleted{\unsure{XI}} \textcircled{$9^\circ$} \uline{The subgroups}
$\mathcal M_{\rho,\sigma}(0)$,
\deleted{$\mathcal M_{\rho,\sigma}(\tfrac12),\mathcal M_{\rho,\sigma}(-\unsure{\overline J|\overline\jmath})$.}
\inserted{$\mathcal M_{\rho,\sigma}[0]$, and the group}
$S\mathcal M_{\rho,\sigma}$, the isomorphism
\diagram{\begin{tikzcd}[ampersand replacement=\&,column sep=large]
S\mathcal M_{\rho,\sigma} \arrow[r,"\sim"] \& \Sigma^+_{\rho,\sigma}
\end{tikzcd}}

The group $\mathcal M_{\rho,\sigma}$ is an extension of $\mathcal T_{\rho,\sigma}$ by
$\unsure{\widetilde G|\widetilde{\mathcal C}}^+_{\rho,\sigma}\simeq SG$, and the three subgroups which we are going
to make explicit will be \deleted{sections}\inserted{\unsure{subgroups|subgroups}}, \deleted{partial} with the
grain of salt that for the third, it will be a
partial section only, over the
subgroup \inserted{$\mathcal T'_{\rho,\sigma}$} of index $2$ in $\mathcal T_{\rho,\sigma}$, kernel of
$\varepsilon_3$.

$[\![$Let us determine in a general way the sections
of $\mathcal M_{\rho,\sigma}$ over $\mathcal T_{\rho,\sigma}$. Since we have

\noindent (1)\quad
\diagram{\begin{tikzcd}[ampersand replacement=\&,column sep=large,row sep=small]
R_{\rho,\sigma} \arrow[r] \arrow[d,phantom,"{\unsure{\wr|\simeq}}" description] \& \mathcal T_{\rho,\sigma} \quad \text{epimorphism with kernel }\mu_2\\
S'T_\rho\times SN'_\sigma\times G_m \&
\end{tikzcd}}

\margin{in the end this won't be of much use\ldots}

it comes to the same to give oneself a section,
or: a homomorphism

\noindent (2)\quad
\diagram{\begin{tikzcd}[ampersand replacement=\&,column sep=large]
R_{\rho,\sigma} \arrow[r,"{\deleted{\unsure{\Xi|\zeta}}\inserted{\psi}}"] \& \mathcal M_{\rho,\sigma}
\end{tikzcd}}

which is \deleted{trivial}
\[
\text{(3)}\qquad
\left\{\begin{aligned}
\text{a)}\quad &\text{trivial on }\mu_2\\
\text{b)}\quad &\deleted{\unsure{\Xi|\Sigma}(x)\unsure{\equiv|=}x\ \bigl(\unsure{\widetilde G|\widetilde{\mathcal C}}^+_{\rho,\sigma}\bigr)\quad\text{identically}.}\quad\inserted{\text{``lifts'' (1)}}
\end{aligned}\right.
\]

The homomorphisms $\unsure{\Xi|\Sigma}$ satisfying (3) correspond
to the splittings of the \uline{trivial} extension of
$R_{\rho,\sigma}$ by $SG$ (\uline{product} extension), hence
simply to the \uline{homomorphisms}

\noindent (4)\quad
\diagram{\begin{tikzcd}[ampersand replacement=\&,column sep=large]
R_{\rho,\sigma} \arrow[r,"{\deleted{\unsure{\Xi|\zeta}_{\unsure{0}}}\inserted{\psi^0}}"] \& SG
\end{tikzcd}},

\mspage{462}

\scanpage{464}

by associating to $\psi^0\deleted{\unsure{\xi_0}}$ the homomorphism
\[
(5)\qquad x\longmapsto \deleted{\unsure{\xi}}\psi(x)
=x\deleted{\unsure{\xi}}\psi^0(x)
=\deleted{\unsure{\xi(x)x}}\psi^0(x).x
\]
Thus condition a takes the form
\[
(6)\qquad \deleted{\unsure{\xi}}\psi^0_{\unsure{0}}(x)=x
\quad\text{on }\mu_2=R_{\rho,\sigma}\cap SG.
\]

\margin{\deleted{\unsure{$\rho$}}}
The $\psi,\psi^0$ corresponding to the \deleted{\unsure{``three''}} sections
$\mathcal M_{\rho,\sigma}^{(0)}$, $\mathcal M_{\rho,\sigma}^{(\frac12)}$,
$\mathcal M_{\rho,\sigma}^{(-\overline J)}$ will be denoted resp.\textsuperscript{ly}
\[
(\psi_{\unsure{0}},\psi^0_0),\ (\psi_{1/2},\psi^0_{1/2}),\
(\psi_{-\overline J},\psi^0_{-\overline J}).\qquad\big]\!\big]
\]

1°) \uline{Case of} $\mathcal M_{\rho,\sigma}^{(0)}$. With the notations
of page 529, it \unsure{will correspond} \deleted{\unsure{$\underline U$ and the}}
\inserted{\unsure{$\underline u$}} of $\mathcal M_{\rho,\sigma}$ satisfying one of the
\uline{equivalent} conditions
\[
(7)\qquad f=1,\quad\inserted{\underline U=\underline u,}\quad
\underline U\in\Gamma T_0,\quad \alpha_0=\alpha,\quad\beta_0=\beta
\]
so we have
\[
(8)\qquad \mathcal M_{\rho,\sigma}^{(0)}=\Theta^{-1}(T_0).
\]
\inserted{(in terms of the description of $\mathcal M_{\rho,\sigma}$ on page 526' and following)}

Let us check that this is a \unsure{subgroup-section}.
We have
\[
\overbrace{\mathcal M_{\rho,\sigma}^{(0)}\cap\widetilde{\unsure{G|\mathcal C}}_{\rho,\sigma}^{+}}
^{\inserted{(\Theta|SG)^{-1}(T_0)}}=\{1\}
\]
\margin{\unsure{and}}
for $\Theta|(\widetilde{\unsure{G|\mathcal C}}_{\rho,\sigma}^{+}=SG)$ is the canonical
inclusion $SG\hookrightarrow G$, now $SG\cap T_0=\{1\}$.
It remains to check that
\diagram{\begin{tikzcd}[ampersand replacement=\&]
\mathcal M_{\rho,\sigma}^{(0)} \arrow[r] \& \mathcal T_{\rho,\sigma}
\end{tikzcd}}
(which is a monomorphism by \unsure{$(a)$|$(\alpha)$}) is an epimorphism. But in

\mspage{463}

\scanpage{465}
\mspage{535}

terms of the canonical \unsure{realization}
\[
(9)\qquad\mathcal M_{\rho,\sigma}\simeq
\underbrace{R_{\rho,\sigma}\times SG}_{\widetilde{\mathcal M}_{\rho,\sigma}}
/\underset{\text{diagonal}}{\mu_2}
\]
and of the homomorphism
\margin{\diagram{\begin{tikzcd}[ampersand replacement=\&]
\kappa:R_{\rho,\sigma}\arrow[r]\&\mathbf G_m
\end{tikzcd}}
the canonical homomorphism \unsure{making it possible} to express $\Theta|R_{\rho,\sigma}$}
\diagram{\begin{tikzcd}[ampersand replacement=\&]
(10)\qquad i_{T_0}:\mathbf G_m\arrow[r,"\sim"]\&T_0\arrow[r,hook]\&G
\end{tikzcd}}
\unsure{isomorphism|homomorphism} inverse to the isomorphism
\diagram{\begin{tikzcd}[ampersand replacement=\&]
\delta_{T_0}=\det|T_0:\quad T_0\arrow[r,"\sim"]\&\mathbf G_m,
\end{tikzcd}}
we have, for an element
\diagram{\begin{tikzcd}[ampersand replacement=\&,column sep=1em,row sep=.2em]
(\underline r,\&U_0)\\
\inserted{R_{\rho,\sigma}}\arrow[u]\&\inserted{SG}\arrow[u]
\end{tikzcd}}
of $\Gamma\widetilde{\mathcal M}_{\rho,\sigma}$
\[
\Theta(\underline r,U_0)=\kappa(\underline r)U_0,\quad
\text{determinant }\kappa(r)^2
\]
to say that this is $\unsure{\in U}$ therefore means we have
\[
\kappa(\underline r)U_0=i_{T_0}(\kappa(\underline r)^2)
\quad\text{i.e. }\deleted{\unsure{U_0=}}
\]
\[
(11)\qquad U_0=\kappa(\underline r)^{-1}i_{T_0}(\kappa(\underline r)^2)
\]
so the inverse image $\widetilde{\mathcal M}_{\rho,\sigma}^{(0)}$ of
$\mathcal M_{\rho,\sigma}^{(0)}$ in $\widetilde{\mathcal M}_{\rho,\sigma}$ consists of the
\[
(12)\qquad
\bigl(\underline r,\kappa(\underline r)^{-1}i_{T_0}(\kappa(\underline r))^2\bigr)
\in\Gamma\widetilde{\mathcal M}_{\rho,\sigma}
\]
in other words this corresponds to a \unsure{group} section, with
\[
(13)\qquad\psi^0_0(\underline r)=\kappa(\underline r)^{-1}i_{T_0}(\kappa(\underline r)^2).
\]
As it should be, we do indeed have here a \uline{homomorphism}
\diagram{\begin{tikzcd}[ampersand replacement=\&]
\psi^0_0:R_{\rho,\sigma}\arrow[r]\&SG
\end{tikzcd}}
and its value on $\mu_2$ is the identity,
for $\kappa(\underline r)^2=1$ and we have
\[
\psi^0_0(r)=\kappa(\underline r)^{-1}=
\begin{gathered}\kappa(\underline r)\\\Vert\\\underline r\end{gathered}
\]

\uline{Remark}\quad\deleted{Before \unsure{qui}}
To define this section, it sufficed

\mspage{464}

\scanpage{466}

(from the point of view of pure calculation) to give oneself an arbitrary homomorphism, denoted
$\unsure{i'_{T_0}|i_{T_0}}$, from $\mathbf G_m$ to $SG$, and to set (13).

Before leaving $\mathcal M_{\rho,\sigma}^{(0)}$, I am going to introduce the
slightly larger group (of dimension 4, instead of 3)
\[
(14)\qquad\mathcal M_{\rho,\sigma}^{[0]}
\overset{\mathrm{dfn}}{=}\Theta^{-1}(B'_0)
\]
where, let us recall, $B'_0$ is defined as the $\frac12$ direct product
\[
(15)\qquad B'_0=T_0.L_0
\]
This time, $U=U_0\kappa(r)\inserted{\in\Gamma B'_0}$ is not determined
by knowledge alone of its determinant\inserted{\unsure{$2$}}
$\underline\kappa(r)^2$, which determines it only modulo $L_0$.
This corresponds to the fact that
$\underset{\inserted{=B'_0\cap SG}}{
\mathcal M_{\rho,\sigma}^{[0]}\cap
\overbrace{\widetilde{\unsure{G|\mathcal C}}_{\rho,\sigma}^{+}}^{\inserted{SG}}}$
is no longer reduced to $(1)$, but to $L_0$, which therefore appears as
an invariant subgroup of $\mathcal M_{\rho,\sigma}^{[0]}$, and we obviously \unsure{will have|have}
(taking into account that $\mathcal M_{\rho,\sigma}^{(0)}\subset\mathcal M_{\rho,\sigma}^{[0]}$ and
\diagram{\begin{tikzcd}[ampersand replacement=\&]
\mathcal M_{\rho,\sigma}^{(0)}\arrow[r,"\sim"]\&\mathcal T_{\rho,\sigma}
\end{tikzcd}}) \unsure{by|that}
\[
(16)\qquad\begin{aligned}
\mathcal M_{\rho,\sigma}^{[0]}&=\mathcal M_{\rho,\sigma}^{(0)}.L_0\\
&\qquad\text{product }\frac12\text{ direct}
\end{aligned}
\]
This is again a construction which makes sense as soon as we have a $T_0$ as above,
\deleted{\illegible{2}} and a subgroup \deleted{\uline{\unsure{affine}}} $L_0$ of $SG$,

\mspage{465}

\scanpage{467}
\mspage{536}

\unsure{invariant} under $T_0$. \deleted{\unsure{We}}

The operation of $\mathcal M_{\rho,\sigma}^{[0]}$ on $SG$ takes place
through $\Theta$, only recently regarded as an epimorphism
\diagram{\begin{tikzcd}[ampersand replacement=\&]
\mathcal M_{\rho,\sigma}^{[0]}
\arrow[r,"\Theta|\mathcal M_{\rho,\sigma}^{[0]}"]\&B'_0
\end{tikzcd}};
this operation being specified, we can introduce
\margin{for a more \unsure{transparent} definition, cf. pages 538', 539}
\[
(17)\qquad\deleted{\unsure{\mathcal S}}
S\mathcal M_{\rho,\sigma}\overset{\text{definition}}{=}
\begin{gathered}
\mathcal M_{\rho,\sigma}^{[0]}.SG\\
\text{product }\frac12\text{ direct}
\end{gathered}
\]
in which $SG$ is invariant, and $\mathcal M_{\rho,\sigma}^{[0]}$ is identified
with a subgroup of $S\mathcal M_{\rho,\sigma}$. Since
$\mathcal M_{\rho,\sigma}^{(0)}\subset\mathcal M_{\rho,\sigma}^{[0]}$,
we have a canonical inclusion
\[
(18)\qquad
\underbrace{\mathcal M_{\rho,\sigma}^{(0)}:SG}_{\unsure{\simeq\mathcal M_{\rho,\sigma}}}
\subset
\underbrace{\mathcal M_{\rho,\sigma}^{[0]}.SG}_{=S\mathcal M_{\rho,\sigma}}
\]
whose image is \deleted{\unsure{denoted} $S\mathcal M_{\rho,\sigma}^{(0)}$}
identified with $\mathcal M_{\rho,\sigma}$.
\deleted{That said, the subgroup $L_0$ of $\mathcal M_{\rho,\sigma}^{[0]}$
cf. (16) \unsure{operates} trivial\textsuperscript{ly} on $SG$,
i.e. $\Theta\inserted{|\mathcal M_{\rho,\sigma}^{[0]}}$,
which is indeed trivial on $L_0$}

\deleted{\unsure{Let us return|Let us resume}} I am beginning \unsure{again} to get lost in it, I am going to hang on
--- for an understanding of the ``geometric'' meaning of
$S\mathcal M_{\rho,\sigma}$, $\mathcal M_{\rho,\sigma}^{[0]}$, to the notations
of page 529, which deserves to be \unsure{made explicit|made explicit} by a

\mspage{466}

\scanpage{468}

\uline{Proposition}\quad Let us consider the homomorphism
\diagram{\begin{tikzcd}[ampersand replacement=\&,column sep=large,row sep=large]
(19)\quad\mathcal M_{\rho,\sigma}\arrow[r]
\&G\times T_\rho\times N'_\sigma\times\mathbf G_m\times
\unsure{\mu_2^2|\mu_2}
\deleted{\unsure{(\mathbf G_m^3\times\mathbf G_m\times\mu_2)}}\\
x\arrow[r,mapsto]
\&(\Theta_G(x),\Theta_\rho(x),\Theta_\sigma(x),\chi(x),\varepsilon_2(x),\varepsilon_3(x))
\end{tikzcd}}
This is a monomorphism, whose image
consists of the \deleted{triples}\inserted{\unsure{systems}}
$(\underline U,\underline a,\underline b,\inserted{\mu,\varepsilon,\varepsilon'})$
satisfying
\[
\deleted{\unsure{\mu=|\text{for}}\,
\unsure{\det_G|\delta_G}(\underline U)
=\delta_\rho(\underline b)=\varepsilon\det}
\]
\[
(20)\qquad\left\{\begin{aligned}
\mu&=\delta_G(\underline U)=\delta_\sigma(\underline b)
=\varepsilon'\delta_\rho(\underline a)\\
\varepsilon&=\varepsilon_\sigma(\underline b)\\[1ex]
\end{aligned}\right.
\]

\uline{Corollary}. The morphism.
\diagram{\begin{tikzcd}[ampersand replacement=\&,column sep=large,row sep=large]
{}\arrow[d,no head,"(21)"']\&\mathcal M_{\rho,\sigma}\arrow[r]
\&G\times T_\rho\times N'_\sigma\\
{}\&x\arrow[r,mapsto]\&\Theta_G(x),\Theta_\rho(x),\Theta_\sigma(x)
\end{tikzcd}}
is a monomorphism, the image consists of the triples
$(\underline U,\underline a,\underline b)$ satisfying
\margin{\uline{N.B.} we have
$\left\{\begin{aligned}
\varepsilon&=\varepsilon_\sigma(\underline b)\\
\varepsilon'&=\delta_\rho(\underline a)/\delta_\sigma(\underline b)\\
\mu&=\delta_\sigma(\underline b)=\delta_G(\underline U)
\end{aligned}\right.$}
\[
(22)\qquad\left\{\begin{aligned}
\delta_G(\underline U)&=\delta_\sigma(\underline b),\\
[\delta_\rho(\underline a)/\delta_\sigma(\underline b)]^2&=1.
\end{aligned}\right.
\]

\uline{Proof of the proposition.}\quad We have
\deleted{$\unsure{\operatorname{Ker}(\Theta_G)|\mu_G}$}
\[
\operatorname{Ker}(\Theta_G,\deleted{\unsure{\chi|\xi_2}},\varepsilon_2,\varepsilon_3)
=SR_{\rho,\sigma}^{\unsure{\circ|0}}\simeq ST_\rho\times ST_\sigma
\]
and $\Theta_\rho,\Theta_\sigma$ \deleted{\unsure{thus|of the}}
on $SR_{\rho,\sigma}^{\unsure{\circ|0}}$ are identified with the two projections,
hence \deleted{$\unsure{\Theta}$}
\[
\operatorname{Ker}\bigl((\Theta_\rho,\Theta_\sigma)|SR_{\rho,\sigma}^{\unsure{\circ|0}}\bigr)
=\{1\},\quad\text{hence}
\]
\[
\operatorname{Ker}(\Theta_G,\chi,\varepsilon_2,\varepsilon_3,\Theta_\rho,\Theta_\sigma)
=(1),
\]

\mspage{467}

\scanpage{469}
\mspage{537}

hence a monomorphism. For the epimorphism, we are reduced to the equivalent form of the Corollary. --- We may assume that we have a $\lambda$ such that
\[
\lambda^2=\mu\overset{\mathrm{dfn}}{=}\det\underline U
\]
we want to write
\[
\underline U=\Theta_{\widetilde{\mathcal M}}
\left(\underset{\underset{R_{\rho\sigma}}{\rotatebox{90}{$\in$}}}{r},
\underset{\underset{SG}{\rotatebox{90}{$\in$}}}{U_0}\right)
=\kappa(r)U_0
\]
\margin{N.B. $\Theta_{\widetilde{\mathcal M}}$ the canonical homomorphism
\diagram{\begin{tikzcd}[ampersand replacement=\&]
\widetilde{\mathcal M}_{\rho,\sigma}\overset{\mathrm{dfn}}{=}R_{\rho,\sigma}\times SG \arrow[r] \& G
\end{tikzcd}}
which I \unsure{had|wrote|write} $\Theta_G$.}
with $U_0\in\Gamma(SG)$, i.e. $\kappa(r)^2=\det\underline U$ i.e. $\kappa(r)^2=\mu$,
we shall take $r$ with \deleted{$\unsure{\kappa}$}
\[
\kappa(r)=\lambda,
\]
\deleted{\unsure{then}} we therefore put
\[
\underline U_0=\lambda^{-1}U,\quad
\underline a_0=\lambda^{-1}\underline a,\quad
\underline b_0=\lambda^{-1}b
\]
\deleted{\unsure{and we are}} \inserted{and we are} reduced to seeing that $\underline U_0,\underline a_0,\underline b_0$ are in the image of $\mathcal M_{\rho,\sigma}$ (in \unsure{being|removing} ``\unsure{corrected ones}'' $(\underline U,\underline a,\underline b)$ by the image \deleted{\unsure{in the}} $G\times T_\rho\times\unsure{T_\sigma}$ of the section $\lambda$ of $G_m\hookrightarrow R_{\rho,\sigma}$). We are reduced to seeing that $(\underline a_0,\underline b_0)$ is in the image of
\diagram{\begin{tikzcd}[ampersand replacement=\&]
R_{\rho,\sigma}\arrow[r] \& \deleted{\unsure{T}}T_\rho\times N'_\sigma.
\end{tikzcd}}
But now $\underline a_0$ is in $S'T_\rho$, $\underline b_0$ in $SN'_\sigma$, and precisely $R_{\rho,\sigma}\simeq S'T_\rho\times SN'_\sigma$, \unsure{we have won}\ldots

\uline{Remark}. We cannot replace $G$ by $PG$,
i.e.
\diagram{\begin{tikzcd}[ampersand replacement=\&]
\mathcal M_{\rho,\sigma}\arrow[r] \& PG\times T_\rho\times N'_\sigma
\end{tikzcd}}
(where we can still add $G_m\times\mu_2\times\mu_2$ if we want!) is \uline{not} a monomorphism --- its kernel is precisely the subgroup $\mu_2$ of $\widetilde{\unsure{G|\mathcal C}}_{\rho,\sigma}^{+}=SG$.

\mspage{468}

\scanpage{470}

Let us consider, now that the quantities $\underline U,\underline a,\underline b$ $(\inserted{and}\ \mu,\varepsilon,\varepsilon')$ of page~529 (152) are understood, the quantities (153) of loc. cit. $\underline u,f,\alpha,\beta,\alpha_0,\beta_0$ ($g,h\ \xi,\eta$ do not appear there, besides\ldots) They are determined in terms of the $\underline U,\underline a,\underline b$, as soon as we know either $\underline u$ or $f$, by the formulas (154), $\underline u$ and $f$ being related by the single relation
\[
\underline U=f^{-1}\underline u \tag{23}
\]
with in addition on $f$ the condition
\[
f\in\Gamma\underbrace{SG}_{\widetilde{\unsure{G|\mathcal C}}_{\rho,\sigma}^{+}}^{\unsure{\sim}},\tag{24}
\]
plus a condition on $u$ which, in loc. cit., was written
\[
\underline u\in\Gamma T_0,\tag{25}
\]
and was suggested by the calculations of \S46 $\overline{\underline{\mathrm V}}$ (taken up again later in \S48 $\overline{\underline{\mathrm{II}}}$, $\overline{\underline{\mathrm{III}}}$ etc). This condition ensures the \uline{uniqueness} of $f,\underline u$ \unsure{satisfying|satisfying} (23) --- that is, from the $\mathrm p^{\mathrm t}$ of view of calculation, its main virtue! --- but in certain contexts it is \unsure{preferable} to replace it by the more \unsure{broad} condition
\[
\underline u\in\Gamma B'_0
\qquad\text{where}\quad
\begin{gathered}
B'_0=T_0.L_0.\\
\text{(product }\tfrac12\text{ direct)}
\end{gathered}\tag{26}
\]

\mspage{469}

\scanpage{471}
\mspage{538}

We no longer have \deleted{any \unsure{result|results}} \uline{uniqueness} for writing $\underline U$ in the form (23) --- on the other hand we have an \uline{existence} theorem in a simple form:

If we \uuline{start} with $(\alpha,\beta,f)$\deleted{$\unsure{\in}$}, sections of $G$ such that
\[
(\det\alpha)^2=\det\beta=\det f=1,\tag{27}
\]
for \deleted{us to be able to write} \inserted{this triple} and as a \unsure{result|consequence} the quantities deduced from it: \deleted{$\alpha_0,\beta_0,$}
\deleted{$\xi,\eta,g,h$}
\[
\left\{
\begin{aligned}
\alpha_0&=f^{-1}\alpha,&\quad \beta_0&=f^{-1}\unsure{\varepsilon}\beta\\
\xi&=\alpha\rho(\alpha)^{-1},&\quad \eta'&=\beta\sigma(\beta)^{-1}\\
g&=\alpha_0\rho(\alpha_0)^{-1},&\quad h&=\beta_0\sigma(\beta_0)^{-1}
\end{aligned}
\right.\tag{28}
\]
to come from a system $(x,\deleted{\unsure{f}}\inserted{\underline u})$,
with $x\in\Gamma\mathcal M_{\rho,\sigma}$,
$\deleted{\unsure{f}}\inserted{\underline u}\in\Gamma\deleted{\unsure{SG|T_0}}\inserted{B'_0}$ such that
\[
f\overset{\mathrm{dfn}}{=}\underline u\,\underline U^{-1}
\deleted{\unsure{\underline u|f}}
\in\Gamma\deleted{\unsure{B'_0}}\inserted{SG},
\quad\text{i.e.}\quad\det\underline u=\det\underline U,
\tag{29}
\]
by the formulas
\[
\alpha=\underline u\,\underline a^{-1},\quad
\beta=\underline u\,\underline b^{-1},\quad
f=\underline u\,\underline U^{-1},\tag{30}
\]
it is necessary and sufficient that we have the relation (which does not involve $f$, only $\alpha,\beta$):
\[
\alpha\rho(\alpha)^{-1}=\varepsilon\beta\sigma(\beta)^{-1}\varepsilon_1(\mu-1)
\qquad
\begin{gathered}
\text{for : }\varepsilon\in\Gamma\mu_2\\
\mu\in\Gamma G_m\\
\text{suitable.}
\end{gathered}\tag{31}
\]

\uline{(N.B.} $\unsure{\varepsilon|\varepsilon_1},\mu$ are uniquely determined (\unsure{thus} in terms of $\alpha\ \beta$).
\inserted{Then the pair $(x,\underline u)$ is unique. For $f=1$, (i.e. $\underline u=\underline U$),}
We can say \deleted{that} also that we have a $1$-$1$ correspondence between the \inserted{the set \deleted{\illegible{1}} sections \unsure{of the|of the}} \unsure{subgroups|subgroup} $\mathcal M_{\rho,\sigma}[0]$ (not $\mathcal M_{\rho,\sigma}(0)$!) of $\mathcal M_{\rho,\sigma}$, and the set

\mspage{470}

\scanpage{472}

of pairs $\alpha,\beta$ satisfying the two relations
\[
(\det\alpha)^2=\det\beta=1\qquad\text{of (27)}
\]
\uline{and} the loop relation. For this, we have \deleted{\unsure{besides}} no use for the group $S\mathcal M_{\rho,\sigma}$.

But returning \unsure{to the case} of arbitrary $f$, to interpret the \unsure{algebraic} \unsure{sheaf} on the triples $(\alpha,\beta,f)$, \inserted{satisfying (27)(31),} \inserted{with the structure of} a strange group on it \inserted{\unsure{via} cocycle relation,} (to which I shall have to return in \inserted{more} \unsure{geometric} terms\ldots) and because the three quantities $\alpha,\beta,f$ mix intimately --- there is reason to introduce $S\mathcal M_{\rho,\sigma}$ --- which will be isomorphic \unsure{to the|to} scheme (in groups) of the common solutions of (27)(31).

From this point of view, $S\mathcal M_{\rho,\sigma}$ is derived from $\unsure{\rho}$ as \deleted{\unsure{described by triples}} \inserted{scheme}
\deleted{$(x,\underline u,f)$ \unsure{of}.}
\[
\begin{aligned}
S\mathcal M_{\rho,\sigma}&=\text{scheme of triples }(x,\underline u,f)\\
&\text{with }\left\{
\begin{aligned}
x\ \deleted{\unsure{\in}}&\text{ in }\mathcal M_{\rho,\sigma},\\
\deleted{\unsure{f}}\inserted{\underline u}&\text{ in }B'_0\\
f&\text{ in }SG
\end{aligned}
\right.
\qquad\left\{
\begin{aligned}
\underline U&=f^{-1}\underline u\\
\Vert\,\text{definition}&\\
\Theta_G(x)&
\end{aligned}\right.
\end{aligned}\tag{32}
\]

Thus
\diagram{\begin{tikzcd}[ampersand replacement=\&,column sep=large,row sep=small]
{(33)\quad S\mathcal M_{\rho,\sigma}}
\arrow[r,hook]
\arrow[d,phantom,"{\unsure{\rotatebox{90}{$\simeq$}}}"]
\& {\mathcal M_{\rho,\sigma}\times B'_0\quad\unsure{\leftarrow}}\\
{\left\{x,\underline u\ \middle|\ \unsure{\chi|\gamma}(x)=
\underbrace{\unsure{\delta_G}(\underline u)}_{\begin{gathered}\Vert\,\mathrm{dfn}\\\det\underline u\end{gathered}}\right\}} \&
\end{tikzcd}}
or again, more simply

\mspage{471}
